\documentclass[11pt,a4paper]{article}
\usepackage[margin=25mm]{geometry}
\usepackage{amsmath,amssymb,amsthm,hyperref}
\newtheorem{theorem}{Theorem}
\newtheorem{lemma}[theorem]{Lemma}
\newcommand{\gen}[1]{\langle#1\rangle}
\title{Enumeration by stabilizers of initial strings}
\author{Human Compression project: F006 recovery}
\date{Working proof; separate review required}
\begin{document}\maketitle
Let $G$ be a specified finite group generated by involutions, and let $J$
be its set of nonidentity involutions. A string marking is an ordered tuple
$(t_1,\ldots,t_m)$ with $[t_i,t_j]=1$ for $|i-j|>1$.
We use the usual intersection property (IP): subgroups on two subsets of
the marks intersect in the subgroup on their common subset. A C-marking
generates its stated group and has IP. Disconnected strings are included.
The standard end criterion of McMullen--Schulte, Proposition 2E16, is a
retained input: two end deletions with IP and the correct common-middle
intersection suffice for a string to have IP.

\begin{lemma}[Appending one mark]\label{append}
Suppose $t=(t_1,\ldots,t_m)$ has IP and $x\in J\setminus\gen t$
commutes with $t_1,\ldots,t_{m-1}$. There is an exact test for IP of
$(t,x)$ using only subgroup generation and intersection. Recursively test
the suffix $(t_2,\ldots,t_m,x)$, and require
\[
 \gen{t_1,\ldots,t_m}\cap\gen{t_2,\ldots,t_m,x}
       =\gen{t_2,\ldots,t_m}.
\]
The recursion ends with tuples of length at most two, which pass.
\end{lemma}
\begin{proof}
Every prefix encountered during this recursion is a submarking of the old
C-marking and therefore has IP. Distinct nonidentity involutions form a
rank-two C-marking: their order-two cyclic subgroups have trivial
intersection. The hypothesis $x\notin\gen t$ ensures distinctness in every
suffix. Induct on the length, applying the end criterion at each step.
Conversely, every suffix and every displayed intersection is required by
IP. The recursion is therefore both necessary and sufficient. Its premise
that the old prefix has IP is essential; the same abbreviated recursion is
not asserted to test arbitrary tuples without that premise.
\end{proof}

Put $A=\operatorname{Aut}(G)$, acting on $J$. This action is faithful:
an automorphism fixing every involution fixes a generating set of $G$.
For each accepted tuple $t$, let $A_t$ be its pointwise stabilizer in $A$.
The following procedure starts with the empty tuple and trivial subgroup.
\begin{enumerate}
\item If $\gen t=G$, record $t$ and do not extend it.
\item Otherwise form every $x\in J\setminus\gen t$ commuting with
all marks except possibly the last. Take one representative from each
$A_t$-orbit on this candidate set.
\item For each representative passing Lemma~\ref{append}, visit $(t,x)$
with generated subgroup $\gen{t,x}$ and stabilizer $(A_t)_x$.
\end{enumerate}
The candidate set is $A_t$-invariant, since its defining membership and
commutation conditions are invariant under automorphisms fixing $t$.

\begin{theorem}\label{complete}
The procedure terminates and records exactly one representative of each
$\operatorname{Aut}(G)$-orbit of generating string C-markings of $G$.
Every possible rank is covered by searching through
$\lfloor\log_2|G|\rfloor$.
\end{theorem}
\begin{proof}
Each appended mark lies outside the previous subgroup, so the subgroup
order grows by an integer index at least two. A visited tuple of length
$m$ consequently generates a subgroup of order at least $2^m$.
This proves the depth bound and termination. A tuple already generating
$G$ has no independent extension. Such pruning loses no C-marking,
because IP implies that no mark lies in the subgroup on the other marks.

Induct on tuple length. The empty tuple has one orbit. Suppose the visited
prefixes represent precisely the orbits of C-markings of length $m$.
Every C-marking of length $m+1$ has one of these prefixes after an
automorphism of $G$. Its final mark is a candidate and passes the exact
test. An element of the prefix stabilizer moves it to the chosen orbit
representative, so the extended marking is visited. Conversely two
extensions can be automorphic only if their prefixes are. Once their
prefix representatives agree, that automorphism lies in $A_t$; distinct
chosen candidate orbits therefore cannot merge. This proves coverage and
uniqueness, both for partial markings and for generating leaves.
\end{proof}

\begin{lemma}[Testing possible completions]
For a visited prefix $t=(t_1,\ldots,t_m)$ put
\[
 P_m=J\cap\bigcap_{i=1}^{m-1}C_G(t_i),\qquad
 K_m=\gen{t_1,\ldots,t_m,P_m}.
\]
An empty intersection means $P_m=J$. If $K_m\ne G$, the whole branch
can be discarded without losing any generating marking. The sets $P_m$
can be maintained by intersection with cached sets $J\cap C_G(t_m)$.
\end{lemma}
\begin{proof}
Every later mark in a string extension has index at least $m+1$ and
therefore commutes with each of $t_1,\ldots,t_{m-1}$. It belongs to
$P_m$, so the completed generated group is contained in $K_m$.
If that group is proper, no completion generates $G$.
For an appended mark $x$, the required pool is precisely
$P_{m+1}=P_m\cap C_G(t_m)$, independently of $x$. Caching the full
commuting set changes neither this equality nor the subgroup test.
These pools are invariant under the pointwise prefix stabilizer.
Thus the pruning preserves the generating automorphism orbits of
Theorem~\ref{complete}. It intentionally changes the counts of visited
partial markings: those counts are no longer a census of all partial
C-markings. The subgroup test uses every element of the pool, not merely
representatives of its stabilizer orbits.
\end{proof}

\begin{lemma}[Duality counts]
Let $N_m$ be the number of recorded markings of rank $m$, and let $S_m$
count those for which an automorphism sends the tuple to its reverse.
There are $(N_m+S_m)/2$ classes up to isomorphism and duality, of which
$S_m$ are self-dual. The same formula holds after restricting to tuples
whose adjacent products all have order greater than two.
\end{lemma}
\begin{proof}
Reversal induces an involution on the finite set of ordered automorphism
orbits. Its fixed points are exactly the $S_m$ specified orbits, and its
other orbits have size two. The connected-string condition is preserved
by automorphisms and reversal. Finally, the correspondence between regular
polytopes and string C-groups identifies an isomorphism of two markings on
the same group with an automorphism sending their distinguished generators
in order. Reversal is precisely duality in that correspondence.
\end{proof}

\paragraph{Implementation and retained inputs.}
The complete GAP program implements these statements literally. One version
forms $J$ from all group elements; the general version unions all conjugacy
classes whose representatives have order two. It checks that their union
has the sum of their class sizes. It computes the full automorphism group
and its faithful permutation action on $J$, checks faithfulness by exact
orders, and uses point stabilizers and candidate orbits. The recursive test
is called only on extensions of already accepted prefixes. Complete leaf
tuples, automorphism generators, involutions and counts at every rank are
retained, including zero counts above the last surviving prefix.

The groups are supplied by exact reflection permutations on finite root
sets. The classical finite Coxeter presentations and orders identify these
models: the simple reflections satisfy the specified relations, and their
generated permutation groups have the classical orders. The complete root
constructor, permutations, group computations and run logs are required
finite evidence, not optional validation. GAP's exact group operations,
complete conjugacy and automorphism algorithms, and their actual native
execution are explicit software premises. This is not a Lean or independently
kernel-checked computation. A timeout is never a classification certificate.
The variant with cached commuting sets and the completion-subgroup test
uses the additional proved lemma above. Its visited-prefix counts are
reported with their pruning counts; only complete generating leaves and
their duality classes are used in the source comparison.

Chen--Rowley, \emph{Abstract Regular Polytopes of Finite Irreducible Coxeter
Groups}, EJC 32(4) (2025), P4.18, supplies the target Table 1. The stabilizer
enumeration is the established method described by Leemans--Mixer,
\emph{Algorithms for classifying regular polytopes with a fixed automorphism
group}, Contributions to Discrete Mathematics 7(2) (2012), 105--118,
Section 5, there attributed to Hartley--Hulpke,
\emph{Polytopes derived from sporadic simple groups}, Contributions to
Discrete Mathematics 5(2) (2010), 106--118. The proof here supplies the
precise extension test and completion-subgroup pruning for our executable
version. It uses ordinary group actions and the retained end criterion;
no priority for these general methods is asserted. A separate
source comparison must distinguish connected and disconnected markings and
record every completed exceptional row. Until those executions and their
review are complete, this note claims no whole-source recovery or numerical
compression factor.
\end{document}
